\chapter{Urinary retention}
\index{Urinary retention}

\section*{Definition and Overview}
Urinary retention is defined as the inability to voluntarily and completely empty the bladder. This condition represents a significant disruption in the normal physiological process of micturition, which relies on a complex, highly coordinated interplay between the central nervous system, the peripheral autonomic pathways, the detrusor muscle of the bladder wall, the bladder neck, and the urethral sphincters. Urinary retention is clinically categorized into two distinct forms: acute urinary retention (AUR) and chronic urinary retention (CUR).

Acute urinary retention is characterized by a sudden, painful, and complete inability to urinate. It is a common urological emergency that requires immediate intervention to decompress the bladder, relieve severe discomfort, and prevent acute kidney injury or bladder rupture. Conversely, chronic urinary retention is a long-standing, progressive condition characterized by a persistent, abnormally high volume of residual urine in the bladder after voiding (referred to as post-void residual volume). Unlike its acute counterpart, chronic retention is often painless or associated with minimal discomfort because the bladder wall gradually distends over months or years, accommodating larger volumes of urine without triggering acute pain receptors. However, this insidious distension can lead to silent renal damage, bladder wall remodeling, and recurrent systemic infections.

The clinical significance of urinary retention cannot be overstated. It is a major cause of morbidity, particularly in the aging male population, and poses a direct threat to renal function. When the bladder fails to empty, intravesical pressure increases and is transmitted retrogradely through the ureters to the kidneys. This back-pressure can lead to bilateral hydronephrosis and post-renal acute kidney injury. Furthermore, stagnant urine serves as a breeding ground for pathogens, predisposing patients to recurrent urinary tract infections, bladder calculi, and potentially fatal urosepsis. Consequently, timely recognition, precise diagnostic evaluation, and appropriate therapeutic intervention are critical to preserving renal function and improving patient outcomes.

\section{*{Epidemiology}}
The epidemiology of urinary retention reveals significant variation across different age groups, biological sexes, and clinical contexts. Acute urinary retention is overwhelmingly a disease of older men, primarily driven by the age-related enlargement of the prostate gland, known as benign prostatic hyperplasia (BPH). Epidemiological studies indicate that the incidence of AUR in men rises exponentially with age. For men in their 60s, the incidence of AUR is estimated at approximately 3 per 1,000 man-years. This rate increases to approximately 10\% for men in their 70s over a five-year period, and approaches nearly 30\% for men in their 80s and older. BPH is implicated in up to 70\% of all cases of AUR in men. Other male-specific factors include prostate cancer and urethral stricture disease.

In stark contrast, acute urinary retention is an exceedingly rare event in women. The estimated annual incidence of AUR in women is approximately 3 per 100,000 individuals. When it does occur in females, it is typically associated with pelvic organ prolapse, gynecological malignancies, pelvic masses, or profound neurological dysfunction.

The epidemiology of chronic urinary retention is more difficult to define precisely, as many individuals remain asymptomatic for extended periods. However, CUR also demonstrates a strong male predominance, mirroring the prevalence of BPH. In younger populations, urinary retention is rare and is almost exclusively associated with congenital anomalies, neurogenic bladder dysfunction (such as that resulting from spina bifida, spinal cord trauma, or demyelinating diseases like multiple sclerosis), or urethral trauma. The socio-economic burden of urinary retention is substantial, encompassing the costs of emergency department visits, catheterization supplies, surgical procedures, and the management of secondary complications such as urinary tract infections and chronic renal failure.

\section{*{Aetiology and Pathophysiology}}
The etiology of urinary retention is multifactorial and can be broadly categorized into obstructive, neurogenic, myogenic, pharmacological, infectious/inflammatory, and psychogenic causes. The underlying pathophysiology involves either a mechanical blockage of urine flow or a functional failure of the detrusor muscle or nervous system coordination.
\begin{itemize}
    \item \textbf{Outflow Obstruction}: This is the most common mechanism leading to urinary retention. In males, the physical obstruction is typically located at the level of the prostatic urethra or bladder neck. Benign prostatic hyperplasia (BPH) causes progressive compression of the urethral lumen due to stromal and epithelial hyperplasia. Prostate cancer can similarly cause mechanical obstruction. In both sexes, urethral strictures, bladder neck contractures, urethral valves, or foreign bodies can physically impede urine flow. In females, pelvic organ prolapse (such as a severe cystocele or uterine prolapse) can cause anatomical distortion and kinking of the urethra. Additionally, external compression from pelvic masses, such as uterine fibroids or ovarian cysts, can obstruct the bladder outlet.
    \item \textbf{Neurologic Dysfunction (Neurogenic Bladder)}: Micturition is regulated by complex neural circuits involving the cerebral cortex, pons, spinal cord (specifically sacral segments S2--S4), and peripheral nerves. Any lesion or disease process that disrupts these pathways can lead to urinary retention. In detrusor-sphincter dyssynergia (DSD), commonly seen in patients with suprasacral spinal cord injuries or multiple sclerosis, there is a lack of coordination between detrusor contraction and external urethral sphincter relaxation, resulting in functional obstruction. Lesions affecting the sacral cord or pelvic nerves (e.g., from cauda equina syndrome, diabetic neuropathy, or pelvic surgery) lead to an underactive or detrusor muscle (flaccid bladder) due to the loss of parasympathetic motor innervation.
    \item \textbf{Detrusor Underactivity (Myogenic Bladder)}: The detrusor muscle itself can lose its contractile strength over time. This is often the result of chronic overdistension. When a patient repeatedly ignores the urge to urinate or when moderate outlet obstruction is left untreated, the detrusor muscle fibers are stretched beyond their physiological limit, leading to ischemic changes, fibrosis, and loss of intercellular gap junctions. Aging also contributes to myogenic changes, including a reduction in detrusor muscle density and altered response to acetylcholine, resulting in detrusor underactivity (hypocontractility).
    \item \textbf{Pharmacological Agents}: A wide variety of medications can precipitate or worsen urinary retention by interfering with autonomic control of micturition. Anticholinergic drugs (such as atropine, oxybutynin, and certain tricyclic antidepressants) block the muscarinic receptors (primarily M3) on detrusor muscle cells, preventing parasympathetic-mediated contraction. Sympathomimetics (such as pseudoephedrine) stimulate alpha-adrenergic receptors at the bladder neck and prostate, increasing outflow resistance. Opioids (such as morphine) act on central and peripheral opioid receptors to inhibit detrusor contractility, reduce bladder sensation, and increase sphincter tone. Antihistamines, calcium channel blockers, and NSAIDs can also contribute to bladder underactivity.
    \item \textbf{Infection and Inflammation}: Acute inflammatory processes in the lower genitourinary tract can lead to urinary retention. Conditions such as acute cystitis, prostatitis, urethritis, and vulvovaginitis can cause localized tissue edema and severe pain during urination. The pain can trigger a reflex spasm of the external urethral sphincter and voluntary or involuntary inhibition of detrusor contraction. Genital herpes infections (particularly primary outbreaks) can cause urinary retention through a combination of local pain and inflammation of the sacral nerve roots (Elsberg syndrome).
    \item \textbf{Postoperative and Traumatic Causes}: Urinary retention is a common postoperative complication, especially following pelvic, spinal, or anorectal surgeries. General and spinal anesthesia temporarily suppress detrusor reflex arcs. Furthermore, intravenous fluid administration during surgery can lead to rapid bladder overdistension while the patient is still anesthetized. Pelvic trauma, including pelvic fractures, can cause direct injury to the urethra or pelvic nerves, leading to retention.
\end{itemize}

\section*{Clinical Features}
The clinical presentation of urinary retention depends heavily on whether the condition is acute or chronic, as well as the underlying etiology.
\begin{itemize}
    \item \textbf{Acute Urinary Retention (AUR)}: The presentation of AUR is sudden, dramatic, and highly distressing. Patients experience a complete and painful inability to pass urine despite a strong, urgent desire to do so. The primary symptom is severe, agonizing pain in the suprapubic and lower abdominal region. On physical examination, a tense, tender, and dull-to-percuss mass is palpable in the suprapubic area, representing the severely distended bladder (often containing 500 mL to over 1000 mL of urine). Patients are typically anxious, diaphoretic, and may exhibit autonomic responses such as hypertension and tachycardia.
    \item \textbf{Chronic Urinary Retention (CUR)}: In contrast to the acute form, CUR develops insidiously and is often asymptomatic or minimally painful, as the bladder wall slowly stretches over time to accommodate large volumes of urine (sometimes exceeding 1500 mL). The clinical features are primarily those of bladder dysfunction and incomplete emptying. Symptoms include hesitancy (difficulty initiating the urinary stream), a weak or interrupted stream, straining to urinate, and a persistent sensation of incomplete bladder emptying. Patients often experience significant urinary frequency (the need to urinate often) and nocturia (waking up multiple times at night to urinate) due to the reduced functional capacity of the bladder.
    \item \textbf{Overflow Incontinence}: This is a hallmark feature of chronic urinary retention. As the bladder reaches its physical limit of distension, intravesical pressure exceeds the outlet resistance, leading to the involuntary leakage of small amounts of urine, often exacerbated by coughing, sneezing, or physical exertion. This is frequently mistaken by patients as simple urinary incontinence, when in fact it represents severe retention.
    \item \textbf{Recurrent Infections and Hematuria}: Stagnant urine in the bladder promotes bacterial growth, leading to recurrent urinary tract infections (UTIs) characterized by dysuria, cloudy urine, and malodour. Chronic mucosal irritation and high intravesical pressure can also cause micro- or macro-hematuria (blood in the urine).
\end{itemize}

\section*{Differential Diagnosis}
It is essential to differentiate urinary retention from other conditions that present with low or absent urine output, or with similar abdominal symptoms.
\begin{itemize}
    \item \textbf{Anuria or Oliguria}: This is the true lack of urine production by the kidneys, secondary to acute kidney injury (AKI), chronic kidney disease (CKD), or severe dehydration, rather than a failure of bladder emptying. Patients with anuria will have an empty bladder, which can be easily confirmed via bladder ultrasound, whereas patients with urinary retention have a full, distended bladder.
    \item \textbf{Benign Prostatic Hyperplasia (BPH)}: Although BPH is a primary cause of retention, it must be differentiated from other causes of bladder outlet obstruction such as prostate cancer, urethral stricture, or bladder neck contracture.
    \item \textbf{Urethral Stricture}: Scar tissue in the urethra, resulting from prior trauma, infection (e.g., gonococcal urethritis), or instrumentation, can mimic the obstructive symptoms of BPH but typically occurs in younger men and requires endoscopic evaluation.
    \item \textbf{Neurogenic Bladder}: Dysfunction resulting from neurological diseases (e.g., multiple sclerosis, spinal cord injury, Parkinson's disease) must be differentiated from mechanical obstruction, as the management strategies differ significantly.
    \item \textbf{Pelvic Organ Prolapse}: In women, severe cystocele, rectocele, or uterine prolapse can physically obstruct the urethra and must be distinguished from primary detrusor underactivity or urethral strictures.
    \item \textbf{Urinary Tract Infection (UTI)}: Acute cystitis or urethritis can cause severe frequency, urgency, and hesitancy. While these can trigger reflex retention, they can also present similarly without significant retention, requiring urinalysis and culture for differentiation.
    \item \textbf{Fecal Impaction}: Severe constipation with a large mass of hard stool in the rectum can mechanically compress the bladder neck and urethra, particularly in elderly, bedridden, or neurologically impaired patients, mimicking primary urinary tract pathology.
    \item \textbf{Acute Abdomen}: The severe abdominal pain associated with acute urinary retention can sometimes be confused with other causes of an acute abdomen, such as appendicitis, diverticulitis, or a ruptured abdominal aortic aneurysm, particularly if the bladder is not palpated or if the patient is unable to provide a clear history.
\end{itemize}

\section*{When to Seek Medical Care}
Understanding when to seek immediate medical attention is crucial for preventing serious complications associated with urinary retention.
\begin{itemize}
    \item \textbf{Emergency Care}: The sudden, complete, and painful inability to urinate is a medical emergency. Delaying treatment can lead to bladder rupture, permanent detrusor muscle damage, or acute kidney failure. If you or someone else experiences this symptom, you must seek emergency medical care immediately. Go to the nearest hospital emergency department or call your local emergency services (e.g., \textbf{local emergency services} in many countries, \textbf{911} in the United States, or \textbf{999} in the United Kingdom).
    \item \textbf{Urgent Symptoms}: You should also seek urgent medical attention if you experience urinary retention accompanied by any of the following systemic symptoms: high fever, chills, severe lower back or flank pain (suggesting pyelonephritis or urosepsis), nausea, vomiting, or progressive leg weakness and numbness in the groin or buttocks area (known as ``saddle anesthesia,'' which indicates cauda equina syndrome, a neurological emergency).
    \item \textbf{Routine Care}: For gradual, chronic symptoms such as a progressively weakening urine stream, difficulty starting urination, frequent urination (especially at night), a persistent feeling that the bladder is not empty, or involuntary leaking of urine (overflow incontinence), you should schedule an appointment with a general practitioner or a urologist for a comprehensive evaluation.
\end{itemize}

\section*{Diagnosis and Investigations}
A systematic approach is required to diagnose urinary retention, determine its severity, identify the underlying cause, and evaluate for complications.
\begin{itemize}
    \item \textbf{Physical Examination}: The initial step is a thorough physical exam. The clinician will palpate and percuss the lower abdomen; a distended bladder is typically felt as a smooth, firm, oval-shaped mass arising from the pelvis, which is dull to percussion. A digital rectal examination (DRE) is essential in male patients to evaluate the size, shape, and consistency of the prostate gland, and to assess anal sphincter tone. In female patients, a pelvic examination is performed to check for pelvic organ prolapse or pelvic masses. A focused neurological examination is also critical, assessing perineal sensation (sacral dermatomes), deep tendon reflexes in the lower extremities, and voluntary anal sphincter contraction.
    \item \textbf{Post-Void Residual (PVR) Measurement}: This is the key diagnostic test for urinary retention. It measures the volume of urine remaining in the bladder immediately after the patient attempts to empty it. PVR is most commonly and non-invasively measured using a portable bedside bladder ultrasound scanner. Alternatively, a quick catheterization can be performed to drain and measure the residual urine. A PVR of less than 50 mL is normal; 50--100 mL is acceptable in older adults. A PVR greater than 100 mL is generally considered abnormal, and a volume exceeding 200 mL strongly suggests urinary retention.
    \item \textbf{Laboratory Investigations}:
    \begin{itemize}
        \item \textbf{Urinalysis and Urine Culture}: Performed to detect the presence of white blood cells, nitrites, or bacteria, indicating a urinary tract infection.
        \item \textbf{Serum Creatinine and Blood Urea Nitrogen (BUN)}: Vital to evaluate renal function. Elevated levels indicate renal impairment, which may be due to post-renal obstruction (obstructive nephropathy).
        \item \textbf{Serum Electrolytes}: Monitored, especially if renal function is impaired, to detect life-threatening abnormalities like hyperkalemia.
        \item \textbf{Prostate-Specific Antigen (PSA)}: May be ordered for men to screen for prostate cancer, though it is important to note that acute urinary retention and bladder catheterization can cause transient, significant elevations in PSA levels.
    \end{itemize}
    \item \textbf{Imaging Studies}:
    \begin{itemize}
        \item \textbf{Renal and Bladder Ultrasound}: Recommended to evaluate the upper urinary tract for hydronephrosis (dilation of the renal pelvis and calyces) and to assess the bladder wall thickness, bladder diverticula, or bladder stones.
        \item \textbf{CT or MRI Scan}: Used if a pelvic or abdominal mass is suspected of causing extrinsic compression on the urinary tract.
    \end{itemize}
    \item \textbf{Urodynamic Studies}: These physiological tests evaluate how well the bladder and urethra store and release urine. They are particularly useful in complex cases, chronic retention, or when neurogenic dysfunction is suspected.
    \begin{itemize}
        \item \textbf{Uroflowmetry}: Measures the rate of urine flow during voiding. A low peak flow rate suggests obstruction or detrusor weakness.
        \item \textbf{Cystometry (Cystometrogram)}: Measures the relationship between bladder pressure and volume during bladder filling, assessing bladder capacity, compliance, and the presence of involuntary detrusor contractions.
        \item \textbf{Pressure-Flow Studies}: Simultaneously measure bladder pressure and flow rate during voiding, helping to differentiate between bladder outlet obstruction and detrusor underactivity.
    \end{itemize}
    \item \textbf{Cystoscopy}: An outpatient endoscopic procedure where a urologist inserts a small, flexible camera (cystoscope) through the urethra into the bladder. This allows direct visualization of the urethral lumen, prostatic urethra, bladder neck, and bladder mucosa to identify strictures, prostatic encroachment, tumors, or stones.
\end{itemize}

\section*{Management}
The management of urinary retention is divided into immediate treatment to relieve the retention and long-term treatment to address the underlying cause.
\begin{itemize}
    \item \textbf{Immediate Decompression}: The primary goal in acute urinary retention is to drain the bladder immediately.
    \begin{itemize}
        \item \textbf{Urethral Catheterization}: The most common method. A flexible tube (Foley catheter) is inserted through the urethra into the bladder, and a balloon at the tip is inflated to hold it in place. The urine is drained into a collection bag.
        \item \textbf{Suprapubic Catheterization}: If a urethral catheter cannot be passed (e.g., due to a severe urethral stricture or acute prostatitis), or if urethral catheterization is contraindicated (e.g., suspected urethral trauma), a catheter is inserted directly into the bladder through a small incision in the lower abdominal wall, just above the pubic bone.
    \end{itemize}
    \item \textbf{Pharmacological Management}:
    \begin{itemize}
        \item \textbf{Alpha-1 Adrenergic Blockers}: Medications such as tamsulosin, alfuzosin, silodosin, and doxazosin relax the smooth muscle fibers in the bladder neck and the prostatic urethra. This reduces resistance to urine flow. They are typically initiated immediately after catheterization to improve the chances of a successful Trial Without Catheter (TWOC).
        \item \textbf{5-Alpha Reductase Inhibitors}: Medications like finasteride and dutasteride block the conversion of testosterone to dihydrotestosterone, leading to a gradual reduction in prostate volume (about 20--25\% over 6 to 12 months). They are used for long-term management of BPH to prevent recurrent retention and the need for surgery.
        \item \textbf{Anticholinergics and Beta-3 Agonists}: While anticholinergics are contraindicated in untreated retention, beta-3 agonists (like mirabegron) or low-dose anticholinergics may be used cautiously in patients with mixed symptoms of retention and overactive bladder, provided their PVR is closely monitored.
    \end{itemize}
    \item \textbf{Clean Intermittent Catheterization (CIC)}: For patients with chronic urinary retention due to neurogenic bladder dysfunction or detrusor underactivity, self-catheterization is the gold standard. The patient is taught to insert a clean, single-use catheter into their urethra to drain the bladder 4 to 6 times a day. This avoids the complications associated with long-term indwelling catheters (such as severe infections and urethral erosion).
    \item \textbf{Surgical and Interventional Procedures}: When conservative and pharmacological treatments fail, surgical intervention is often required, particularly for anatomical obstruction.
    \begin{itemize}
        \item \textbf{Transurethral Resection of the Prostate (TURP)}: The historical gold standard surgical treatment for BPH. A resectoscope is inserted through the urethra, and electrical energy is used to shave away obstructing prostate tissue.
        \item \textbf{Laser Prostatectomy}: Techniques such as Holmium Laser Enucleation of the Prostate (HoLEP) or Photoselective Vaporization of the Prostate (PVP) use laser energy to remove or vaporize obstructing tissue. HoLEP is highly effective for very large prostates and has a lower risk of bleeding.
        \item \textbf{Minimally Invasive Surgical Therapies (MISTs)}: Options like the prostatic urethral lift (UroLift) or water vapor therapy (Rezum) provide symptom relief with lower risks of sexual side effects compared to TURP.
        \item \textbf{Urethral Dilation and Urethroplasty}: For urethral strictures, the stricture can be dilated using specialized instruments or surgically reconstructed (urethroplasty).
        \item \textbf{Sacral Neuromodulation (InterStim)}: An implantable device that stimulates the sacral nerves, sometimes used for patients with idiopathic urinary retention or non-obstructive urinary retention.
    \end{itemize}
    \item \textbf{Lifestyle and Behavioral Therapy}:
    \begin{itemize}
        \item \textbf{Double Voiding}: Instructing the patient to urinate, wait a few minutes, and then try to urinate again to ensure complete bladder emptying.
        \item \textbf{Timed Voiding}: Urinating on a set schedule (e.g., every 3 to 4 hours) rather than waiting for the urge, which helps prevent bladder overdistension.
        \item \textbf{Pelvic Floor Muscle Training}: Kegel exercises can help improve sphincter control and coordination.
    \end{itemize}
\end{itemize}

\section*{Complications}
Untreated or poorly managed urinary retention can lead to severe and potentially permanent complications.
\begin{itemize}
    \item \textbf{Urinary Tract Infections (UTIs) and Pyelonephritis}: Stagnant urine in the bladder acts as an incubator for pathogenic bacteria. Recurrent UTIs are common, and the infection can ascend to the kidneys, causing acute pyelonephritis, which can lead to permanent renal scarring.
    \item \textbf{Urosepsis}: This is a life-threatening systemic inflammatory response to a urinary tract infection. Bacteria enter the bloodstream from the infected urinary tract, leading to septic shock, multi-organ failure, and death if not treated urgently with intravenous antibiotics and supportive care.
    \item \textbf{Obstructive Nephropathy and Renal Failure}: Chronic high intravesical pressure is transmitted back through the ureters to the renal pelvis, causing hydroureter and hydronephrosis. Over time, this pressure damages the renal parenchyma, leading to chronic kidney disease (CKD) and potentially end-stage renal disease (ESRD) requiring dialysis.
    \item \textbf{Bladder Remodeling and Detrusor Muscle Failure}: Chronic distension causes hypertrophy of the detrusor muscle, collagen deposition, and fibrosis. This leads to a loss of bladder compliance and contractility, culminating in a flaccid, non-functional bladder (myogenic bladder failure) that cannot empty without catheterization.
    \item \textbf{Bladder Calculi and Diverticula}: Urinary stasis and chronic infection promote the precipitation of mineral crystals, leading to the formation of bladder stones. High intravesical pressure can also cause the bladder mucosa to herniate through the detrusor muscle fibers, forming bladder diverticula (outpouchings), which further worsen stasis.
    \item \textbf{Post-Obstructive Diuresis (POD)}: Following the rapid decompression of a chronically distended bladder, some patients experience a dramatic, pathological increase in urine output (often exceeding 200 mL per hour). This occurs due to the loss of the renal medullary concentration gradient and transient tubular dysfunction. POD can lead to severe dehydration, hypovolemia, and life-threatening electrolyte imbalances (such as hypokalemia and hyponatremia) requiring close monitoring and intravenous fluid replacement.
\end{itemize}

\section*{Prognosis and Outcomes}
The prognosis for individuals with urinary retention varies widely depending on the underlying cause, the duration of the retention, the patient's overall health, and the speed of medical intervention. For patients presenting with acute urinary retention (AUR) secondary to benign prostatic hyperplasia (BPH), the immediate prognosis is excellent with catheterization. Following catheter insertion and the initiation of alpha-blockers, a Trial Without Catheter (TWOC) is successful in approximately 50--60\% of patients. However, a significant proportion of these patients will experience a recurrence of AUR within one year, and many will ultimately require surgical intervention (such as TURP or HoLEP). The long-term prognosis after successful surgery for BPH-related obstruction is outstanding, with most patients experiencing a dramatic improvement in urinary symptoms, flow rate, and overall quality of life.

In cases where urinary retention is caused by acute infections or medication side effects, the prognosis is highly favorable once the infection is treated or the offending medication is discontinued. For patients with chronic urinary retention (CUR) or retention due to neurogenic bladder dysfunction (e.g., spinal cord injury or advanced multiple sclerosis), the prognosis is more guarded. While these conditions are often not ``curable'' in the traditional sense, they can be successfully managed long-term using clean intermittent catheterization (CIC). Patients who adhere to a strict CIC regimen and receive regular urological follow-up can maintain excellent renal function and lead active, fulfilling lives. The primary factors associated with a poorer prognosis include delayed diagnosis (leading to irreversible detrusor muscle fibrosis or advanced kidney disease), persistent urosepsis, and poor compliance with self-care or catheterization regimens.

\section*{Prevention and Self-Care}
While not all cases of urinary retention can be prevented, several proactive measures and lifestyle choices can significantly reduce the risk of developing this condition or prevent its recurrence.
\begin{itemize}
    \item \textbf{Medication Caution}: Be highly vigilant when taking over-the-counter (OTC) medications. Decongestants containing pseudoephedrine or phenylephrine, and antihistamines (such as diphenhydramine) commonly found in cold, allergy, and sleep aids, have strong anticholinergic or sympathomimetic properties that can precipitate acute urinary retention, especially in men with BPH. Always consult a pharmacist or doctor before starting these medications.
    \item \textbf{Healthy Bladder Habits}: Avoid delaying urination. Empty your bladder as soon as you feel the urge, as holding urine for prolonged periods can stretch and weaken the detrusor muscle over time. Practice double voiding if you feel you are not emptying completely: after urinating, remain on the toilet, relax for a minute, and attempt to void again.
    \item \textbf{Diet and Fluid Management}: Maintain adequate, consistent hydration throughout the day, but avoid drinking large volumes of fluid in a short period. Limit consumption of bladder irritants such as caffeine, alcohol, and highly spicy foods, which can exacerbate urinary urgency and hesitancy.
    \item \textbf{Preventing Constipation}: A full rectum can press against the bladder neck and urethra, causing obstruction. Prevent constipation by eating a high-fiber diet rich in fruits, vegetables, and whole grains, staying physically active, and drinking sufficient water.
    \item \textbf{Pelvic Floor Health}: Perform pelvic floor muscle exercises (Kegel exercises) regularly to strengthen the muscles that support the bladder and regulate urination. This is particularly beneficial for women with mild pelvic organ prolapse and men recovering from prostate surgery.
    \item \textbf{Regular Medical Monitoring}: Men over the age of 50 should undergo regular prostate health screenings and discuss any changes in their urinary patterns (such as hesitancy, weak stream, or frequent nighttime urination) with their general practitioner to address BPH before it leads to retention.
\end{itemize}

\section*{Sources and Further Reading}
For reliable, evidence-based information regarding urinary retention and general urological health, please refer to the following resources:
\begin{itemize}
    \item National Institute of Diabetes and Digestive and Kidney Diseases (NIDDK). \textit{Urinary Retention}. Available at: \texttt{https://www.niddk.nih.gov/health-information/urologic-diseases/urinary-retention} (Accessed: August 2026).
    \item American Urological Association (AUA). \textit{Guidelines on Benign Prostatic Hyperplasia (BPH) and Lower Urinary Tract Symptoms}. Available at: \texttt{https://www.auanet.org/guidelines-and-quality/guidelines/benign-prostatic-hyperplasia-(bph)-guideline} (Accessed: August 2026).
    \item Urology Care Foundation. \textit{What is Urinary Retention?}. Available at: \texttt{https://www.urologyhealth.org/urology-a-z/u/urinary-retention} (Accessed: August 2026).
    \item Healthdirect Australia. \textit{Urinary Retention}. Available at: \texttt{https://www.healthdirect.gov.au/urinary-retention} (Accessed: August 2026).
    \item NHS Choices. \textit{Urinary Retention}. Available at: \texttt{https://www.nhs.uk/conditions/urinary-retention/} (Accessed: August 2026).
    \item European Association of Urology (EAU). \textit{Patient Information on Urinary Retention}. Available at: \texttt{https://patients.uroweb.org/post/urinary-retention/} (Accessed: August 2026).
\end{itemize}